\section*{Chapter \ref{ch:g6:for-loops-turtle} --- Repeating with \texttt{for}: Turtle Drawings}

\begin{solution}{exo:g6:for-loops-turtle:1}
Three lines: \code{0}, \code{1} and \code{2}.
\end{solution}

\begin{solution}{exo:g6:for-loops-turtle:2}
2, 5 and 8: from 2 in steps of 3, stopping before 10.
\end{solution}

\begin{solution}{exo:g6:for-loops-turtle:3}
12 lines: the outer \omterm{def:g4:algorithms-repeat-and-decide:loop}{loop} runs 3 times, and each time the inner \omterm{def:g4:algorithms-repeat-and-decide:loop}{loop} prints 4
lines, from \code{0 0} to \code{2 3}.
\end{solution}

\begin{solution}{exo:g6:for-loops-turtle:4}
Each turn is $360 \div 3 = 120$ degrees.

\omcode{code/grade-6/03-for-loops-turtle/exercises/triangle.py}{1}{6}{An
equilateral triangle.}\label{lst:g6:for-loops-turtle:sol-triangle}
\end{solution}

\begin{solution}{exo:g6:for-loops-turtle:5}
A pentagon needs turns of $360 \div 5 = 72$ degrees; with 60 the five turns
make only 300 degrees, less than a full turn, so the turtle does not come back
to its start (the drawings below). The fix is \code{left(72)}.
\end{solution}

\begin{solution}{exo:g6:for-loops-turtle:6}
Indent the second \code{print}, so that it is in the \omterm{def:g6:for-loops-turtle:for-loop}{block}:

\omcode{code/grade-6/03-for-loops-turtle/exercises/hip.py}{1}{3}{Both lines in
the \omterm{def:g6:for-loops-turtle:for-loop}{block}.}\label{lst:g6:for-loops-turtle:sol-hip}
\end{solution}

\begin{solution}{exo:g6:for-loops-turtle:7}
\omcode{code/grade-6/03-for-loops-turtle/exercises/seven.py}{1}{2}{The 7 times
table.}\label{lst:g6:for-loops-turtle:sol-seven}
\end{solution}

\begin{solution}{exo:g6:for-loops-turtle:8}
Each turn is $360 \div 8 = 45$ degrees.

\omcode{code/grade-6/03-for-loops-turtle/exercises/octagon.py}{1}{6}{A regular
octagon.}\label{lst:g6:for-loops-turtle:sol-octagon}
\end{solution}

\begin{solution}{exo:g6:for-loops-turtle:9}
$360 \div 10 = 36$ degrees: going once around the polygon, the turtle makes one
full turn, $360^\circ$, shared equally between the 10 corners.
\end{solution}

\begin{solution}{exo:g6:for-loops-turtle:10}
After each triangle, the turtle is back at the triangle's first corner, facing
right; one more \code{forward(50)} takes it to the end of the base, where the
next triangle starts.

\omcode{code/grade-6/03-for-loops-turtle/exercises/row_of_triangles.py}{1}{8}{A
row of triangles.}\label{lst:g6:for-loops-turtle:sol-triangles}

\begin{omfigure}
\begin{tikzpicture}
  \foreach \f/\t/\x/\w in {triangle/{triangle.py}/0/3.4, open_polygon/{open\_polygon.py}/3.3/3.4, octagon/{octagon.py}/6.6/3.4, row_of_triangles/{row\_of\_triangles.py}/9.9/5.5} {
    \begin{axis}[at={(\x cm,0)}, width=\w cm, height=3.4cm, axis equal image, hide axis,
        title={\footnotesize\ttfamily\t}, title style={yshift=-1ex}]
      \addplot[very thick, ocBlue!80!black, mark=none] table[header=false, x index=0, y index=1]
        {code/grade-6/03-for-loops-turtle/figures/out/\f.dat};
      \addplot[only marks, mark=*, mark size=1.4pt, ocRed] coordinates {(0,0)};
    \end{axis}
  }
\end{tikzpicture}

{\small The drawings of exercises 4, 5, 8 and 10. With turns of 60 degrees,
the ``pentagon'' stops short of its start.}
\end{omfigure}
\end{solution}

\begin{solution}{exo:g6:for-loops-turtle:11}
$5 \times 144 = 720 = 2 \times 360$: two full turns. After the five sides the
turtle faces right again, as at the start. It is also back at its start: the
five sides are the same length and the five turns the same, so the five points
lie evenly on a circle, and the turtle jumps two points at a time until it has
gone twice around the circle and reached the first point again.
\end{solution}

\begin{solution}{exo:g6:for-loops-turtle:12}
The number of rank $i$, counted from 0, is $1 + 3 \times i$:

\omcode{code/grade-6/03-for-loops-turtle/exercises/sequence.py}{1}{3}{The first
20 numbers.}\label{lst:g6:for-loops-turtle:sol-sequence}

The twentieth is $1 + 3 \times 19 = 58$.
\end{solution}

\begin{solution}{pb:g6:for-loops-turtle:1}
\textbf{1.} $360 \times 1 \div 5 = 72$ degrees: a regular pentagon.

\textbf{2.} $360 \times 2 \div 5 = 144$ degrees: the five-pointed star of
\cref{lst:g6:for-loops-turtle:star}, drawn turning left, so it is its mirror
image.

\textbf{3.} Jumping two places gives the five-pointed star. Jumping three places
gives the same star, drawn the other way round: a jump of 3 forwards is a jump
of 2 backwards. In the same way, $k = 3$ turns by $216^\circ$ to the left, which
is $360 - 216 = 144^\circ$ to the right: the star again.

\begin{omfigure}
\begin{tikzpicture}
  \foreach \k [count=\i from 0] in {1,2,3} {
    \pgfmathsetmacro{\x}{3.2*\i}
    \begin{axis}[at={(\x cm,0)}, width=3.4cm, height=3.4cm, axis equal image, hide axis,
        title={\scriptsize $n = 5$, $k = \k$}, title style={yshift=-1.2ex}]
      \addplot[thick, ocBlue!80!black, mark=none] table[header=false, x index=0, y index=1]
        {code/grade-6/03-for-loops-turtle/figures/out/star-5-\k.dat};
      \addplot[only marks, mark=*, mark size=1.3pt, ocRed] coordinates {(0,0)};
    \end{axis}
  }
\end{tikzpicture}

{\small Five points: the pentagon, and the star drawn in both directions.}
\end{omfigure}

\textbf{4.} $360 \times 2 \div 8 = 90$ degrees: the turtle turns as for a
square, but draws 8 sides, so it goes around the square twice.

\textbf{5.} $360 \times 4 \div 8 = 180$ degrees: the turtle turns right round
after each side, so it walks back and forth on one segment, 4 times each way.

\textbf{6.} $k = 1$, 3, 5 and 7.

\textbf{7.} $k = 5$ turns $225^\circ$ to the left, which is $135^\circ$ to the
right, while $k = 3$ turns $135^\circ$ to the left: the same star, drawn in
mirror image. $k = 7$ turns $315^\circ$ left, that is $45^\circ$ right: the
octagon of $k = 1$, drawn the other way. In general, $k$ and $n - k$ give mirror
drawings.

\textbf{8.} $k = 1$, 5, 7 and 11. For $k = 3$, the turn is $90^\circ$: a square,
gone around three times. For $k = 4$, it is $120^\circ$: a triangle, gone around
four times.

\textbf{9.} Counting the drawings:

\begin{center}
\begin{tabular}{|l|c|c|c|c|c|c|c|c|}
\hline
$n$ & 5 & 6 & 7 & 8 & 9 & 10 & 11 & 12 \\ \hline
\omterm{def:g6:variables-and-input:value}{values} of $k$ & 4 & 2 & 6 & 4 & 6 & 4 & 10 & 4 \\ \hline
\end{tabular}
\end{center}
For example, for $n = 6$ only $k = 1$ and $k = 5$ work (the hexagon); for
$n = 9$, every $k$ except 3 and 6.

\textbf{10.} The divisors of 8 other than 1 are 2, 4 and 8: the $k$ sharing none
of them are the odd ones, 1, 3, 5, 7, as in question 6. The divisors of 12 other
than 1 are 2, 3, 4, 6 and 12: the $k$ from 1 to 11 that are neither even nor
multiples of 3 are 1, 5, 7 and 11, as in question 8.

\textbf{11.} The only divisors of a prime number $n$ are 1 and $n$ itself. A
number $k$ from 1 to $n - 1$ is too small to have $n$ as a divisor, so $k$ and
$n$ share only the divisor 1, and every $k$ works: 6 \omterm{def:g6:variables-and-input:value}{values} for 7, 10 for 11.

\textbf{12.} \textbf{Result.} For 12 points, only $k = 1$, 5, 7 and 11 draw all
12 points in one line: $k = 1$ and 11 the regular 12-sided polygon, $k = 5$
and 7 a 12-pointed star; the other \omterm{def:g6:variables-and-input:value}{values} of $k$ go several times around a
smaller shape.
\end{solution}
